% ============================================================================
%  DRAFT v1 -- Physical Review Letters
%  Emergent altermagnetism and d_xy pairing in the checkerboard Hubbard model
%
%  Every number in this file traces to a file in the repository. Pointers are
%  given in comments beside each claim so the draft can be checked against data
%  rather than against memory. Items still needing input are marked TODO.
% ============================================================================
\documentclass[aps,prl,reprint,superscriptaddress,amsmath,amssymb,floatfix]{revtex4-2}

\usepackage{graphicx}
\usepackage{bm}

\begin{document}

\title{Emergent altermagnetism and $d_{xy}$ pairing in the checkerboard Hubbard model}

% TODO: author list and affiliations to be completed by the authors.
\author{[Author One]}
\affiliation{[Affiliation]}
\author{[Author Two]}
\affiliation{[Affiliation]}

\date{\today}

\begin{abstract}
We study the half-filled Hubbard model on the checkerboard lattice by
constrained-path quantum Monte Carlo. The Hamiltonian is spin independent, so
the two magnetic sublattices are related by a fourfold rotation rather than by a
translation, and any spin splitting of the momentum distribution must be
generated by the interaction. We find a compensated spin-split state above a
finite interaction strength whose momentum-space polarization is odd under
$C_4$ and even under the diagonal mirror at all 126 interacting parameter sets
studied, which identifies the $d_{xy}$ channel. A control calculation at
vanishing lattice anisotropy, where the model reduces to the square lattice and
the splitting is forbidden by symmetry, fixes the statistical floor of the order
parameter at $6.3(2)\times10^{-4}$ per site, two orders of magnitude below the
measured signal. The interaction-induced pairing vertex vanishes identically in
the noninteracting limit and is largest in the $d_{xy}$ channel over the same
range of anisotropy where the magnetic signal is $d_{xy}$. Magnetism and pairing
therefore select the same representation, which distinguishes this state from
the $d_{x^2-y^2}$ pairing reported for the checkerboard lattice in the regime
where the two models coincide.
\end{abstract}

\maketitle

% ---------------------------------------------------------------------------
\section{Introduction}
% ---------------------------------------------------------------------------

Altermagnets combine a vanishing net moment with a spin-split band structure,
a combination long assumed to require either ferromagnetism or spin-orbit
coupling~\cite{Smejkal2022a,Smejkal2022b}. The splitting follows from the
crystal symmetry: the two opposite-spin sublattices are connected by a rotation
rather than by a translation or an inversion, so the spin degeneracy that a
conventional antiferromagnet enforces is absent.

Most microscopic models of this state build the splitting into the Hamiltonian,
usually through a spin-dependent hopping. The splitting is then present already
at zero interaction and the question of whether interactions alone can produce
it is left open. Two recent studies of the Lieb lattice state the gap directly,
noting that unbiased studies establishing an altermagnet as the ground state of
an interacting Hamiltonian are lacking~\cite{Kaushal2025,Duerrnagel2025}. Both
use methods with a controlled range of validity, unrestricted Hartree-Fock with
exact diagonalization on small clusters in the first case and a functional
renormalization group treatment of the weak-coupling limit in the second.

The checkerboard lattice offers a route to the same physics without putting it
in by hand. It is a square lattice with diagonal bonds on alternating
plaquettes, so the two sublattices are exchanged by a fourfold rotation about a
plaquette center but by no translation. The hopping is spin independent, and
the interaction is a uniform on-site Hubbard $U$. Any spin splitting therefore
appears only once magnetic correlations develop, which is the emergent route
rather than the built-in one.

We report constrained-path quantum Monte Carlo (CPQMC) calculations for this
model at half filling on lattices of $8\times8$, $10\times10$, and $12\times12$
sites. Three results follow. The momentum-resolved spin polarization is nonzero
above a finite $U$ and has $d_{xy}$ symmetry at every interacting parameter set
we examined. A null calculation at vanishing anisotropy, where symmetry forbids
the signal, places the statistical floor two orders of magnitude below the
measured values. The interaction-induced pairing vertex is largest in the
$d_{xy}$ channel over the same anisotropy range, so the magnetic and pairing
instabilities lie in the same representation.

% ---------------------------------------------------------------------------
\section{Model and method}
% ---------------------------------------------------------------------------

We consider
\begin{equation}
H = \sum_{ij\sigma} t_{ij}\, c^{\dagger}_{i\sigma} c^{\phantom{\dagger}}_{j\sigma}
    + U \sum_i n_{i\uparrow} n_{i\downarrow},
\label{eq:H}
\end{equation}
with nearest-neighbor hopping $t_{ij} = t_0$ and diagonal hopping
$t_1 + t_2$ along $(1,1)$ and $t_1 - t_2$ along $(1,-1)$ on one sublattice, the
two exchanged on the other. Energies are in units of $|t_0|$, and we set
$t_0 = -1$ and $t_1 = 0.3$ throughout. The anisotropy is parametrized by
$\delta$ through $t_2 = -\delta$, so that $\delta = 0$ restores the plain square
lattice with uniform second-neighbor hopping and both sublattices become
equivalent. All calculations are at half filling.

The hopping in Eq.~(\ref{eq:H}) is spin independent by construction, so
$n_{\uparrow}(\bm{k}) = n_{\downarrow}(\bm{k})$ at $U = 0$ for every $\delta$.

We use CPQMC with a Trotter step $\Delta\tau = 0.01$, a projection length
$\beta = 32$, and $N_w = 1000$ walkers. The constrained-path approximation
introduces a bias controlled by the trial wave function, so we do not describe
these results as exact. Doubling the projection length from $\beta = 32$ to
$\beta = 48$ changes the order parameter by between $0.1\%$ and $4.1\%$ with no
systematic drift, which we take as convergence in $\beta$.
% beta convergence: 32 -> 48 changes 0.1, 1.5, 2.4, 4.1 percent

The order parameter is the momentum-space spin polarization
\begin{equation}
\Delta_{\mathrm{tot}} = \sum_{\bm{k}} \left| n_{\uparrow}(\bm{k}) - n_{\downarrow}(\bm{k}) \right|,
\label{eq:dtot}
\end{equation}
following Ref.~\cite{OurPRB}. The absolute value inside the sum makes
$\Delta_{\mathrm{tot}}$ sensitive to a compensated splitting, for which the
unsigned sum $\sum_{\bm{k}} [n_{\uparrow}(\bm{k}) - n_{\downarrow}(\bm{k})]$
vanishes. We quote $\Delta_{\mathrm{tot}}/N$ so that different lattice sizes can
be compared, and we note that Eq.~(\ref{eq:dtot}) as written in
Ref.~\cite{OurPRB} is a bare sum.

Because $H$ commutes with the total spin, the polarization must be induced by a
symmetry-breaking field. We add a staggered field $h$ that couples with opposite
sign to the two sublattices, which makes the trial state symmetry broken, and we
report the behavior as a function of $h$. The constrained path is defined by the
trial state, so this field biases the walkers toward the broken state, and the
control described below is what establishes that the measured signal is not an
artifact of that bias.

To classify the symmetry of $\Delta n(\bm{k}) = n_{\uparrow}(\bm{k}) - n_{\downarrow}(\bm{k})$
we use two dimensionless ratios,
\begin{align}
A_{\mathrm{odd}} &= \frac{\sum_{\bm{k}} \left| \Delta n(\bm{k}) + \Delta n(R_{90}\bm{k}) \right|}
                         {\sum_{\bm{k}} \left| \Delta n(\bm{k}) \right|},
\label{eq:aodd}\\
M_{\mathrm{odd}} &= \frac{\sum_{\bm{k}} \left| \Delta n(\bm{k}) + \Delta n(M\bm{k}) \right|}
                         {\sum_{\bm{k}} \left| \Delta n(\bm{k}) \right|},
\label{eq:modd}
\end{align}
where $R_{90}$ is the fourfold rotation and $M$ the diagonal mirror
$(k_x,k_y) \to (k_y,k_x)$. A state odd under $C_4$ gives
$A_{\mathrm{odd}} = 0$ and even gives $A_{\mathrm{odd}} = 2$. A $d_{x^2-y^2}$
state gives $M_{\mathrm{odd}} = 0$ and a $d_{xy}$ state gives
$M_{\mathrm{odd}} = 2$. These ratios use only the symmetry of the data, so
unlike a projection onto a single lattice harmonic they are not sensitive to
how the weight is distributed among higher harmonics of the same symmetry.

% ---------------------------------------------------------------------------
\section{Emergent altermagnetic order}
% ---------------------------------------------------------------------------

$\Delta_{\mathrm{tot}}$ vanishes identically at $U = 0$ and becomes finite for
$U > 0$ at every anisotropy studied, reaching $\Delta_{\mathrm{tot}}/N$ between
$2.2\times10^{-3}$ and $9.2\times10^{-2}$ over the interacting parameter sets.
% range from data/fortran_L12_all49.csv, U>0
The splitting is compensated, so the state carries momentum-resolved spin
polarization with no net moment.

The symmetry classification is unambiguous. Across all 126 interacting parameter
sets at $L = 8$, $10$, and $12$ we obtain $A_{\mathrm{odd}}$ between $0.004$ and
$0.263$ and $M_{\mathrm{odd}}$ between $1.24$ and $2.00$.
% data/classify_3L.csv, 126 cells
Restricted to $L = 12$ the ranges are $0.008$ to $0.166$ and $1.54$ to $2.00$.
% data/classify_L12.csv, 42 cells
The polarization is therefore odd under $C_4$ and even under the diagonal
mirror at every parameter set, which is the $d_{xy}$ channel. No parameter set
is classified as $d_{x^2-y^2}$.

% ---------------------------------------------------------------------------
\section{Null control at vanishing anisotropy}
% ---------------------------------------------------------------------------

Equation~(\ref{eq:dtot}) sums absolute values, so statistical noise contributes
with a definite sign and $\Delta_{\mathrm{tot}}$ has a positive floor that does
not vanish when the physics requires it to. Establishing the scale of that floor
requires a calculation in which the signal is forbidden but the sampling is
unchanged.

The $U = 0$ column does not serve this purpose. There the Green function is
obtained without auxiliary fields, so the zero is a matter of arithmetic rather
than of statistics and it tests nothing about the sampling.

At $\delta = 0$ both diagonal bonds are equal on every site, the two sublattices
become equivalent, and the model reduces to the square lattice with uniform
second-neighbor hopping. Its N\'eel sublattices are related by a translation, so
$n_{\uparrow}(\bm{k}) = n_{\downarrow}(\bm{k})$ identically and
$\Delta_{\mathrm{tot}}$ must vanish by symmetry at any $U$. Whatever is measured
there is the statistical floor at production settings.

We find $\Delta_{\mathrm{tot}}/N = 6.3(2)\times10^{-4}$, averaged over
$U = 2$, $3$, $3.5$, $4$, $4.5$, and $5$ at $L = 12$.
% data/delta0_control_L12.csv: 0.000626 +/- 0.000214
At $\delta = 0.3$ the measured signal ranges from $1.47\times10^{-2}$ at
$U = 2$ to $7.13\times10^{-2}$ at $U = 4.5$, which is between $23$ and $114$
times this floor.
% data/fortran_L12_all49.csv, delta=0.3

% ---------------------------------------------------------------------------
\section{Pairing in the same representation}
% ---------------------------------------------------------------------------

We compute the pair-field susceptibility in four channels, on-site $s$,
extended $s$, $d_{x^2-y^2}$, and $d_{xy}$, and separate the interaction-induced
vertex from the uncorrelated part. The vertex is the meaningful measure of
whether the interaction builds pairing, since the full correlator is large and
positive already for free electrons.

Three features stand out. The vertex is zero in all four channels at $U = 0$,
so everything that follows is interaction generated. The on-site $s$ vertex is
negative throughout, as a repulsive $U$ requires. The $d_{xy}$ vertex is the
largest of the four for $\delta \geq 0.3$ at every $U$ studied, and it grows
monotonically with $U$, from $3.347(8)$ at $U = 2$ to $5.952(22)$ at
$U = 4.5$, in both cases peaking at $\delta = 0.4$.
% data/chi_L12_vertex_summary.csv

The leading channel changes with anisotropy. For $\delta \leq 0.2$ the extended
$s$ channel is largest, and for $\delta \geq 0.3$ the $d_{xy}$ channel is. This
crossover falls in the same range of $\delta$ where the slope of
$\Delta_{\mathrm{tot}}$ with respect to $U$ changes sign, so two quantities
measured independently change character together.

The magnetic and pairing instabilities therefore select the same
representation. The polarization is $d_{xy}$ at all 126 parameter sets, and the
pairing vertex is largest in $d_{xy}$ over the anisotropy range where that
classification holds most strongly.

% ---------------------------------------------------------------------------
\section{Robustness and comparison}
% ---------------------------------------------------------------------------

At half filling with periodic boundaries the noninteracting gap of
Eq.~(\ref{eq:H}) is zero at every $L$ and $\delta$ we examined, so the occupied
set is fixed by the trial state rather than by the spectrum. We repeated the
calculation with antiperiodic boundaries, which shift the momentum mesh off the
high-symmetry points and open a gap, at otherwise identical settings.

The sign of $\partial \Delta_{\mathrm{tot}} / \partial U$ agrees between the two
boundary conditions at $\delta = 0.1$, where both are negative, and at
$\delta = 0.5$, where both are positive, with magnitudes matching to within a
few percent despite the different meshes. At $\delta = 0.3$ the two disagree,
and the antiperiodic value is consistent with zero. We read $\delta \approx 0.3$
as the crossover where the slope passes through zero, at which point the sign is
most sensitive to the mesh, and we therefore do not quote a slope at that
anisotropy.
% data/bc_test, analyze_bc_test.py

The closest published work is a determinant quantum Monte Carlo study of the
checkerboard Hubbard model that reports $d$-wave and $d + is$ pairing as the
dominant channels~\cite{Pan2024}. The two Hamiltonians are not the same. That
model keeps one diagonal per sublattice and omits the other, whereas
Eq.~(\ref{eq:H}) keeps both, and the two coincide only where our second diagonal
vanishes, at $\delta = t_1 = 0.3$. We verified that the single-particle spectra
agree there to within $4\times10^{-15}$.

At that point, and at half filling, we evaluated the vertex in the basis of
Ref.~\cite{Pan2024} alongside our own. Within their basis the $s$-wave channel
is the larger of the two at $U = 3$, $4$, and $4.5$. Including $d_{xy}$, which
their basis does not contain, gives a $d_{xy}$ vertex larger than either.
% data/pan, analyze_pan.py
The comparison is restricted, since Ref.~\cite{Pan2024} works at density
$n = 0.9$ and temperature $T/t = 1/6$ while our calculations are at half filling
in the ground state, and under the particle-hole transformation that relates the
two signs of the second-neighbor hopping their density maps to $n = 1.1$. The
statement is that the dominant channel we find lies outside the basis used
there, not that the two calculations disagree.

% ---------------------------------------------------------------------------
\section{Conclusion}
% ---------------------------------------------------------------------------

The half-filled checkerboard Hubbard model supports a compensated spin-split
state that is absent at $U = 0$ and is $d_{xy}$ in symmetry at every parameter
set examined. A control at vanishing anisotropy, where the splitting is
forbidden, places the statistical floor two orders of magnitude below the
signal. The interaction-induced pairing vertex is largest in the same $d_{xy}$
channel, and the anisotropy at which the pairing channel changes coincides with
the anisotropy at which the magnetic response changes sign.

The mechanism uses no spin-dependent hopping and no spin-orbit coupling. It
requires only a lattice whose magnetic sublattices are related by a rotation and
a uniform Hubbard repulsion, which suggests that the same construction should
apply to other lattices with the same sublattice relation. Whether the state
survives to the thermodynamic limit is not settled by the sizes reached here,
and the constrained-path bias remains.

\begin{acknowledgments}
% TODO: funding and computing acknowledgments.
\end{acknowledgments}

\begin{thebibliography}{99}

% TODO: verify all bibliographic details before submission. Entries marked
% INCOMPLETE need the full reference supplied by the authors.

\bibitem{Smejkal2022a}
L.~\v{S}mejkal, J.~Sinova, and T.~Jungwirth,
Phys. Rev. X \textbf{12}, 031042 (2022).

\bibitem{Smejkal2022b}
L.~\v{S}mejkal, J.~Sinova, and T.~Jungwirth,
Phys. Rev. X \textbf{12}, 040501 (2022).

\bibitem{Kaushal2025}
N.~Kaushal and M.~Franz,
Phys. Rev. Lett. \textbf{135}, 156502 (2025).

\bibitem{Duerrnagel2025}
M.~D\"urrnagel \textit{et al.},
Phys. Rev. Lett. \textbf{135}, 036502 (2025);
arXiv:2412.14251.

\bibitem{Pan2024}
Y.~Pan, R.~Ma, C.~Chen, Z.~Jia, and T.~Ma,
arXiv:2409.16523.

\bibitem{OurPRB}
% INCOMPLETE: this is the group's published PRB defining Eq. (3). Full
% reference to be supplied by the authors.
[Authors], Phys. Rev. B \textbf{[vol]}, [page] ([year]).

\end{thebibliography}

\end{document}
